Güvenlik açısından kritik, yaşlanan elektronik donanımın modernizasyonu, tamamen değiştirilmesine kıyasla çoğu zaman daha uygun maliyetli ve çevresel açıdan daha sürdürülebilirdir; ancak bu, özgün sistemle işlevsel eşdeğerliğin gösterilebilmesine ve modernize edilmiş sistemin güncel standartlara göre yeniden ser\-ti\-f\kern0pt i\-ka\-lan\-dı\-rı\-la\-bil\-me\-sine bağlıdır. Önceki araştırmaları sürdüren bu vaka çalışması, söz konusu zorluğu Avrupa kentsel raylı ulaşımında kullanılan, bir tekerlek kayma koruma sistemine ait eski bir test ve tanılama modülü örneği üzerinden ele almaktadır.

Yenileme (retrof\kern0pt it) kapsamında, Intel~MCS-48 tabanlı, 8~bitlik eski bir denetleyici kartı, özgün merkezî işlem birimini (CPU) ve giriş/çıkış (I/O) mantığını emüle eden, flash tabanlı bir \texttt{Actel~A3P1000} Sahada Programlanabilir Kapı Dizisi (FPGA) ile donatılmış; bu yapı, gelişmiş veri kaydı için galvanik olarak yalıtılmış, 32~bitlik bir Arm\textsuperscript{\textregistered} mikrodenetleyici (µC) olan \texttt{STM32F401RET6} ile tamamlanmıştır. Merkezî bir hedef, güvenlik açısından kritik sistemlerin yenilenmesinde daha geniş çapta uygulanmaya elverişli, metodolojik bir test çerçevesi oluşturmaktı. Bu çerçevenin özü, V~modelinden türetilmiş iki katmanlı bir doğrulama stratejisidir: (1) Bileşen ve Arayüz Testleri: Deterministik donanım tanımlama dili (HDL) simülasyonu, JTAG kontrolleri ve donanım zamanlama ölçümleri, mantıksal doğruluğu, donanımın temel bütünlüğünü ve arayüz zamanlamasını (örneğin watchdog zamanlayıcısı, frekans taraması) incelemektedir. (2) Sistem Düzeyi Testleri: Bu amaca özel bir test düzeneğinde~\cite{dolata2023}, YAML ile tanımlanmış test senaryolarıyla yürütülen senaryo tabanlı testler, entegre davranışı (arıza belleği, öz testler, tanılama) incelemektedir. Simülasyon ve entegre senaryo testleri, güvenlik açısından kritik çekirdeğin arıza işlemedeki (arızanın algılanması, kaydedilmesi, gösterilmesi ve okunması) işlevsel eşdeğerliğine işaret etmektedir. Ancak frekans taraması sapmaktadır ve bu sapma henüz açıklığa kavuşmamıştır: Özgün kartla yapılan tek karşılaştırmalı ölçümde emülasyonun taraması yaklaşık 70~Hz daha düşük bir değerde sona ermekte (emülasyonda yaklaşık 850~Hz, özgün kartta yaklaşık 920~Hz) ve tasarım değerlerine göre,~\cite{dolata2023} numaralı kaynağın özgün test kartı için verdiği yaklaşık 180~ms süren frekans değişiminden daha hızlı ilerlemektedir; sapmanın açıklığa kavuşturulması için özgün kartın yeniden ölçülmesi ya da tarama devresinin kontrol edilmesi gerekmektedir. Donanımda yapılan iki değişiklik, modülün davranışını tasarım gereği değiştirmektedir: Sistem konektörünün STORE ve RECALL girişleri artık arıza belleğine etki etmemekte, eklenen bir düşük gerilim algılama devresi ise emüle edilen denetleyiciyi sıfırlayabilmektedir; bunların araç için önem taşıyıp taşımadığı henüz belirsizdir. Kanıtlar henüz eksiksiz değildir: Arşivlenmiş simülasyonlar ve ölçümler 2023 yılına, yani yeniden yapım öncesine aittir; yeniden yapılan cihazda ise çalışmaları devralan bir ekip, test düzeneğinin test paketini arşivlenmiş bir kayıt olmaksızın yürütmüştür ve aralarında öz testin de bulunduğu birkaç senaryo bu cihaz için belgelenmemiştir. Yeniden yapım sırasında belgelenen altı entegrasyon hatası, nasıl bulunduklarına göre sınıflandırılmaktadır; hatalar, iki katmanlı yöntemin planlı bir uygulamasında gün yüzüne çıkmamıştır ve hiçbir yöntem bu hataların tümünü tek başına bulamamıştır. Katmanlar, IEC~61508~Ed.~2, EN~50129:2018, EN~50716:2023 ve IEEE~1012:2024 standartlarıyla bir bütün olarak ve uygunluk değerlendirmesi yapılmaksızın eşleştirilmiştir; EN~50155:2021 gelecekteki çalışmalara bırakılmıştır. Döngüde donanım (HIL) testine yönelik eksiksiz sistemler bu projenin kapsamı dışında kalmış olsa da yaklaşım, ileride yapılacak bir entegrasyonu kolaylaştırmak için HIL uygulamalarından yararlanmaktadır.

Mevcut raylı araç elektroniğinin elden geçirilmesi maliyetleri düşürebilir ve yeni bileşenlerin üretiminden kaynaklanan gömülü emisyonların önüne geçebilir; kart düzeyindeki yenilemeler de aynı gerekçeye dayanır. Var olan f\kern0pt i\-lo\-larda da kullanılması öngörülen ve yürürlükteki EN~15595 standardına göre ser\-ti\-f\kern0pt i\-ka\-lan\-dı\-rıl\-mış uyarlamalı bir tekerlek kayma koruma (WSP) işlevi bir emsal oluşturmaktadır; metodolojinin başka yenileme vakalarında geçerliliğinin sınanması ise gelecekteki çalışmalara bırakılmıştır.
